\section{Bound on IdealToIsogeny}
\label{sec:idealtoisogeny}

We now revisit the integer-size bound for
Algorithm~\ref{alg:IdealToIsogeny}.  The dominant arithmetic in this procedure
comes from the ideal operations used to prepare suitable short representatives.
In the previous analysis, this contribution was bounded using a coarse
worst-case estimate for ideal multiplication.  This was sufficient for the
correctness and asymptotic structure of the algorithm, since the remaining steps
were already bounded by the same dominant term.  However, after refining the
analysis of ideal multiplication, the same improvement can be propagated through
the subroutines used by \textsf{IdealToIsogeny}.

The first step is to update the bound for
Algorithm~\ref{alg:SuitableIdeals}.  This algorithm computes a reduced basis of
the input ideal and then searches for two short elements whose normalized norms
satisfy a Bezout relation modulo a power of two.  The search phase itself does
not introduce the largest integers: as in the previous analysis, the maximum
size occurs during the ideal multiplication required to compute the reduced
basis.  Therefore, replacing the old ideal multiplication estimate by
Lemma~\ref{lem:idealmult-bound} immediately yields a tighter bound for
\textsf{SuitableIdeals}.

\begin{algorithm}
\caption{SuitableIdeals$(I)$.}\label{alg:SuitableIdeals}
  \begin{algorithmic}[1]
    \Statex \textbf{Input:} A left $\mathcal{O}_0$-ideal $I$, a bound $d=\Theta(\sqrt{p})$.
    \Statex \textbf{Output:} $\beta_1,\beta_2\in I$ and $u,v\in\mathbb{N}^*$ and $f\le e$ such that $\gcd(u q_I(\beta_1),v q_I(\beta_2))=1$ and $u q_I(\beta_1)+v q_I(\beta_2)=2^f$.
    \State Compute a Minkowski reduced basis $(\alpha_1,\dots,\alpha_4)$ of $I$.
    \State Let $B_j\leftarrow \left\lfloor \frac{1}{4}\sqrt{\frac{d}{q_I(\alpha_j)}}\right\rfloor$ for $j\in[1;4]$.
    \State Sample $x_j\in [-B_j;B_j]^4$ for $j\in[1;4]$ and initialize $L\leftarrow [(x_1,\dots,x_4)]$.
    \For{$(y_1,\dots,y_4)\in [-B_1;B_1]\times\cdots\times[-B_4;B_4]$}
      \For{$(x_1,\dots,x_4)\in L$}
        \State $\beta_1\leftarrow \sum_{j=1}^4 x_j\alpha_j$ and $\beta_2\leftarrow \sum_{j=1}^4 y_j\alpha_j$.
        \State $d_1\leftarrow q_I(\beta_1),\ d_2\leftarrow q_I(\beta_2)$.
        \If{$d_1\equiv 1\pmod{2}$ and $d_2\equiv 1\pmod{2}$ and $\gcd(d_1,d_2)=1$}
          \State $u\leftarrow 2^{e} d_1^{-1}\bmod d_2$.
          \State $v\leftarrow \frac{2^{e}-u d_1}{d_2}$.
          \If{$v>0$}
            \State $e'\leftarrow \min(\nu_2(u),\nu_2(v))$, $u\leftarrow u/2^{e'}$, $v\leftarrow v/2^{e'}$, and $f\leftarrow e-e'$.
            \State \textbf{return} $\beta_1,\beta_2,u,v,f$.
          \Else
            \State Append $(y_1,\dots,y_4)$ to $L$.
          \EndIf
        \EndIf
      \EndFor
    \EndFor
  \end{algorithmic}
\end{algorithm}

The following lemma records the resulting refined bound.  Compared with the
previous accounting, the algorithm is unchanged; only the size analysis is
updated.  The improvement comes from using the sharper bound for the intermediate
ideal product appearing in the computation of the Minkowski reduced basis.

\begin{lemma}
During the execution of Algorithm~\ref{alg:SuitableIdeals}, the maximum size of integers is less than \[\frac{1024p^2}{\pi^4}\log{p} \cdot \nrd(I).\]
\end{lemma}

\begin{proof}
To compute a Minkowski reduced basis, one forms an ideal $\bar{J_t}I$ and its Gram matrix, then computes LLL-reduced basis~\cite{AAA+25}, where $J_t$ is the SQIsign parameter ideal.
By the proof of~\cite[Lemma~6]{KLKL25}, the maximum size of integers occurs by the ideal multiplication in Line~1.
Then, by Lemma~\ref{lem:idealmult-bound}, Lemma~\ref{lem:denombound}, and the fact that $\nrd(\bar{J_t}) < (\log{p})^2$, (See~\cite[Appendix~E]{KLKL25}.) we have the maximum size less than
\[
    \frac{1024p^2}{\pi^4}\log{p} \cdot \nrd(I).
\]
\end{proof}
We next propagate this refined estimate to
Algorithm~\ref{alg:IdealToIsogeny}.

\begin{algorithm}
\caption{IdealToIsogeny($I$)}\label{alg:IdealToIsogeny}
\begin{algorithmic}[1]
\Statex \textbf{Input:} A left $\mathcal{O}_0$-ideal $I$.
\Statex \textbf{Output:} The codomain $E_I$ of $\varphi_I$, $\varphi_I(P_0)$, and $\varphi_I(Q_0)$.
\State $u,v,e,s,t,\beta_1,\beta_2 \gets {\textsf{SuitableIdeals}}(I)$
\State $d_1 \gets \mathrm{nrd}(\beta_1)/\mathrm{nrd}(\overline{J_s}I)$
\State $d_2 \gets \mathrm{nrd}(\beta_2)/\mathrm{nrd}(\overline{J_t}I)$ \hfill {// $ud_1 + vd_2 = 2^e$, with $\gcd(ud_1,vd_2)=1$}
\State $E_u,\varphi_u(P_s), \varphi_u(Q_s) \gets {\textsf{FixedDegreeIsogeny}}(s,u)$ \hfill {// $\varphi_u: E_s \to E_u$ is an isogeny of degree $u$}
\State $E_v,\varphi_v(P_t), \varphi_v(Q_t) \gets {\textsf{FixedDegreeIsogeny}}(t,v)$ \hfill {// $\varphi_v: E_t \to E_v$ is an isogeny of degree $v$}
\State $[P,Q]^T \gets \left[\frac{1}{\mathrm{nrd}(I)\mathrm{nrd}(J_t)}\right] \mathbf{M}_{\beta_1} \mathbf{M}_{\beta_2} [\varphi_v(P_t), \varphi_v(Q_t)]^T$
\State $K_P \gets [2^{f-e}]([d_1]\varphi_u(P_s), P)$
\State $K_Q \gets [2^{f-e}]([d_1]\varphi_u(Q_s), Q)$
\State $E \times E', [(P,P'),(Q,Q')] \gets {\textsf{Isogeny22Chain}}(K_P, K_Q, [(\varphi_u(P_s), 0_{E_v}), (\varphi_u(Q_s), 0_{E_v})])$
\If{$e_{2f}(P,Q) = e_{2f}(P_s, Q_s)^{u^2 d_1}$} \hfill {// We identify the right curve $E_I$ by ensuring it is $d_1$-isogeneous to $E_0$}
    \State $E_I \gets E$, $P_I \gets P$, and $Q_I \gets Q$
\Else
    \State $E_I \gets E'$, $P_I \gets P'$, and $Q_I \gets Q'$
    \State $[P_I,Q_I]^T \gets \left[\frac{1}{ud_1}\right]\mathbf{M}_{\beta_1}[P_I,Q_I]^T$
\EndIf
\State \textbf{return} $E_I, (P_I, Q_I)$
\end{algorithmic}
\end{algorithm}

The improved bound for ideal multiplication, once inserted into \textsf{SuitableIdeals}, also gives the refined bound for the full \textsf{IdealToIsogeny} procedure.

\begin{lemma}
\label{lem:IdealToIsogeny-bound}
During the execution of Algorithm~\ref{alg:IdealToIsogeny}, the maximum size of integers is less than \[\frac{1024p^2}{\pi^4}\log{p} \cdot \nrd(I).\]
\end{lemma}

\begin{proof}
By following the same process of the proof of~\cite[Lemma~7]{KLKL25}, we have the same bound with Algorithm~\ref{alg:SuitableIdeals}.
\end{proof}